\begin{resumo}

    \noindent
        A expansão da internet e dos meios de comunicação ampliou o acesso a conteúdos em diferentes idiomas, como vídeos, artigos e blogs. Nesse contexto, as ferramentas de tradução automática contribuem para a redução das barreiras linguísticas ao permitir a conversão de informações entre diferentes línguas. O presente trabalho compara três modelos de tradução automática neural baseados na arquitetura Transformer e disponibilizados com pesos abertos: MarianMT, M2M-100 e NLLB. Os modelos foram avaliados na tradução do inglês para o português, considerando textos de domínio geral e específico. A qualidade das traduções foi analisada por meio das métricas BLEU, BERTScore e chrF, enquanto a eficiência computacional foi avaliada a partir do uso médio de CPU, do tempo de inferência e do pico de memória RAM. Nas condições experimentais estabelecidas, o MarianMT apresentou os melhores resultados nas três métricas de qualidade, além dos menores valores de uso médio de CPU e tempo de inferência. Embora tenham sido observadas diferenças entre os modelos quanto ao processamento, os valores de pico de memória RAM foram próximos. Também se verificou que os textos específicos resultaram em pontuações inferiores às dos textos gerais nas métricas de qualidade. Portanto, no cenário analisado, o MarianMT apresentou o melhor compromisso entre qualidade de tradução e eficiência computacional, enquanto o pico de memória RAM não permitiu diferenciar significativamente os modelos.
    \noindent

    \textbf{Palavras-chave}: transformer, tradução automática, métricas automáticas, eficiência computacional.

\end{resumo}
